\chapter{基于昇腾 310B 的 Ubuntu
教程}\label{ux9644ux5f55-2ux57faux4e8eux6607ux817e-310b-ux7684-ubuntu-ux6559ux7a0b}

本附录把 Ubuntu 基础操作、远程访问、网络配置和昇腾 310B
板端命令整理为一条可复用的学习路径。它面向三类位置：负责编辑代码和文档的\textbf{开发机}、负责连接网络和发起远程操作的
\textbf{Ubuntu 主机}、真正运行 CANN、PyACL、OM 推理和摄像头程序的
\textbf{310B
板端}。同一个命令在不同位置执行，得到的结果和证明范围不同，因此每一步都要先确认命令在哪台机器、哪个用户、哪个
shell 和哪个软件环境中运行。

本附录的命令来自 Ubuntu
基础教程中的文件结构、终端编辑器、SSH、VNC、NFS、Docker、热点与
WiFi、开机自启动、静态 IP 和无屏幕联网资料，并结合现有 310B
材料重新编排。外部资料中的 ROS 小车、树莓派、RDK X5、Orin
等名称只作为操作背景出现于命令示例的原始路径中，不能据此推断 310B
板的实际设备名、IP 地址、接口编号或运行结果。

\section{1. 教程定位与证据边界}\label{src-appendix2-scope}

\subsection{三类操作位置}\label{ux4e09ux7c7bux64cdux4f5cux4f4dux7f6e}

{ % do not increment counter
\begin{longtable}[]{@{}
  >{\raggedright\arraybackslash}p{(\linewidth - 6\tabcolsep) * \real{0.2500}}
  >{\raggedright\arraybackslash}p{(\linewidth - 6\tabcolsep) * \real{0.2500}}
  >{\raggedright\arraybackslash}p{(\linewidth - 6\tabcolsep) * \real{0.2500}}
  >{\raggedright\arraybackslash}p{(\linewidth - 6\tabcolsep) * \real{0.2500}}@{}}
\toprule\noalign{}
\begin{minipage}[b]{\linewidth}\raggedright
位置
\end{minipage} & \begin{minipage}[b]{\linewidth}\raggedright
用途
\end{minipage} & \begin{minipage}[b]{\linewidth}\raggedright
能做什么
\end{minipage} & \begin{minipage}[b]{\linewidth}\raggedright
不能据此推断什么
\end{minipage} \\
\midrule\noalign{}
\endhead
\bottomrule\noalign{}
\endlastfoot
开发机 & 编辑正文、样例代码和配置文件 & 文本编辑、Python 语法检查、纯
Python 测试、前端构建、静态文件检查 & 不能推断板端
CANN、PyACL、OM、摄像头、音频和 NPU 结果 \\
Ubuntu 主机 & 作为 SSH、VNC、NFS、SCP 和 rsync 的客户端 &
发起远程连接、挂载共享目录、传输文件、运行普通 Ubuntu 命令 &
不能因为客户端命令成功就认定板端推理或硬件测试成功 \\
310B 板端 & 运行目标 Ubuntu 系统与昇腾运行时 & 加载 CANN、执行 ATC、导入
acl、运行 OM、检查设备、服务和日志 &
不能把一次导入成功写成完整模型精度或端到端性能 \\
\end{longtable}
}

开发机和 Ubuntu 主机可以共享普通 Linux
操作，例如查看路径、编辑文本、管理网络和传输文件。板端操作则需要额外的证据控制：\passthrough{\lstinline!source /usr/local/Ascend/ascend-toolkit/set\_env.sh!}、\passthrough{\lstinline!npu-smi info!}、\passthrough{\lstinline!atc!}、\passthrough{\lstinline!python -c 'import acl'!}、OM
加载、V4L2 摄像头访问和真实 NPU 推理都应在 310B 板上执行。

\subsection{板端专属操作}\label{ux677fux7aefux4e13ux5c5eux64cdux4f5c}

以下操作默认只能在已经安装好驱动、CANN 和目标 Python 环境的 310B
板端执行：

\begin{lstlisting}[language=bash]
source /usr/local/Ascend/ascend-toolkit/set_env.sh
command -v atc
python -c 'import acl; print("PyACL import: ok")'
npu-smi info
ls -l /dev/video*
v4l2-ctl --device=/dev/video0 --list-formats-ext
\end{lstlisting}

\passthrough{\lstinline!command -v atc!} 检查实际调用的 ATC
路径，\passthrough{\lstinline!import acl!} 检查当前 Python 是否能够加载
PyACL，\passthrough{\lstinline!npu-smi info!} 读取 NPU
设备摘要。三者属于不同的证据层：ATC 存在不代表转换成功，PyACL
可导入不代表模型能运行，\passthrough{\lstinline!npu-smi!}
能输出也不代表某个 OM 已经通过数值和性能验证。

\subsection{证据状态}\label{ux8bc1ux636eux72b6ux6001}

记录结果时，使用下面五类标签，避免把先验知识写成已经观测到的事实：

{ % do not increment counter
\begin{longtable}[]{@{}ll@{}}
\toprule\noalign{}
标签 & 含义 \\
\midrule\noalign{}
\endhead
\bottomrule\noalign{}
\endlastfoot
\passthrough{\lstinline!documented!} &
官方文档、厂商文档或教程中明确写出的能力 \\
\passthrough{\lstinline!inferred!} &
根据已有资料推断出的可能路径，尚未在目标板上验证 \\
\passthrough{\lstinline!observed-pass!} &
在当前板端、当前环境和当前命令下实际通过 \\
\passthrough{\lstinline!observed-fail!} &
在当前组合下实际失败，并保留了完整错误和命令 \\
\passthrough{\lstinline!untested!} & 尚未执行，不能归入支持或不支持 \\
\end{longtable}
}

一条命令的通过只能证明它对应的最小验证门。语法检查、HTTP 健康检查、ATC
转换、ACL
数值烟测、任务精度、性能测试和界面烟测必须分别记录。没有执行的板端测试不能写成``已验证''，没有采集到的日志不能补写。

\subsection{执行命令前的四个问题}\label{ux6267ux884cux547dux4ee4ux524dux7684ux56dbux4e2aux95eeux9898}

执行任何会改变文件、网络、用户、软件包或设备状态的命令前，先回答：

\begin{enumerate}
\def\labelenumi{\arabic{enumi}.}
\tightlist
\item
  命令在哪台机器、哪个 shell 和哪个 conda 环境中运行？
\item
  输入文件、输出目录、日志和远端路径是否已经解析为明确的绝对路径？
\item
  这一步证明的是语法、路径、网络、转换、数值一致性、性能，还是硬件现象？
\item
  命令是否安装软件、修改权限、覆盖文件、删除数据或暴露网络服务？
\end{enumerate}

四个问题没有答案时，先执行只读检查，不要直接执行写入命令。

\section{2. Ubuntu 与 Linux
基础}\label{ubuntu-ux4e0e-linux-ux57faux7840}

\subsection{终端、用户和提示符}\label{ux7ec8ux7aefux7528ux6237ux548cux63d0ux793aux7b26}

Ubuntu 的图形界面之外，终端是与 Linux 交互的主要入口。提示符
\passthrough{\lstinline!$!}
通常表示普通用户，\passthrough{\lstinline!\#!} 通常表示 root 用户。进入
root 状态以后，所有命令的权限都会提高，误操作的后果也会放大。

\begin{lstlisting}[language=bash]
whoami
id
hostname
pwd
printf 'shell=%s\n' "$SHELL"
\end{lstlisting}

\passthrough{\lstinline!whoami!}
显示当前用户名，\passthrough{\lstinline!id!}
同时显示用户、组和补充组，\passthrough{\lstinline!hostname!}
显示主机名，\passthrough{\lstinline!pwd!}
显示当前目录，\passthrough{\lstinline!$SHELL!} 显示当前 shell
路径。远程登录后先确认这些信息，避免把命令发到另一台同名设备。

退出临时 root 状态或当前 shell 可使用：

\begin{lstlisting}[language=bash]
exit
\end{lstlisting}

也可以按 \passthrough{\lstinline!Ctrl+D!}
发送文件结束符。不要为了省事长期保留 root 终端，更不要在 root
终端中编辑项目文件。

\subsection{第一条环境检查命令}\label{ux7b2cux4e00ux6761ux73afux5883ux68c0ux67e5ux547dux4ee4}

\begin{lstlisting}[language=bash]
uname -a
cat /etc/os-release
date --iso-8601=seconds
ip addr
ip route
\end{lstlisting}

\passthrough{\lstinline!uname -a!}
查看内核和架构，\passthrough{\lstinline!/etc/os-release!} 查看 Ubuntu
版本，\passthrough{\lstinline!date!}
生成带时区含义的时间戳，\passthrough{\lstinline!ip addr!}
查看接口地址，\passthrough{\lstinline!ip route!}
查看默认路由。把这几条命令的输出和后续测试放在同一份记录中，才能说明测试环境。

\subsection{可审计 shell}\label{ux53efux5ba1ux8ba1-shell}

脚本或排障会话可以使用严格模式：

\begin{lstlisting}[language=bash]
set -euo pipefail
pwd
whoami
hostname
printf 'shell=%s\n' "$SHELL"
\end{lstlisting}

\passthrough{\lstinline!set -e!}
在未处理的命令失败后停止脚本，\passthrough{\lstinline!-u!}
让未定义变量直接报错，\passthrough{\lstinline!pipefail!}
保留管道前段的失败状态。某个命令允许失败时，应显式处理返回值，不能为了继续执行而全局关闭严格模式。

交互式排查时可以先只运行环境检查命令。脚本中如果必须临时关闭
\passthrough{\lstinline!-u!}
来加载厂商环境脚本，应在加载完成后恢复原状态：

\begin{lstlisting}[language=bash]
load_cann() {
  if [[ $- == *u* ]]; then
    set +u
    source /usr/local/Ascend/ascend-toolkit/set_env.sh
    set -u
  else
    source /usr/local/Ascend/ascend-toolkit/set_env.sh
  fi
}
\end{lstlisting}

\subsection{命令权限与执行位置}\label{ux547dux4ee4ux6743ux9650ux4e0eux6267ux884cux4f4dux7f6e}

同一句命令在开发机、Ubuntu 客户端和板端的结果可能不同。使用
\passthrough{\lstinline!command -v!} 查看实际可执行文件：

\begin{lstlisting}[language=bash]
command -v python
command -v python3
command -v atc || true
command -v npu-smi || true
\end{lstlisting}

\passthrough{\lstinline!|| true!}
只适合让检查流程继续，不代表后面的命令可以跳过环境问题。若
\passthrough{\lstinline!atc!} 或 \passthrough{\lstinline!npu-smi!}
不存在，应记录当前机器和当前环境，而不是修改代码来掩盖缺失的板端依赖。

\section{3.
文件系统、路径与常用命令}\label{ux6587ux4ef6ux7cfbux7edfux8defux5f84ux4e0eux5e38ux7528ux547dux4ee4}

\subsection{Ubuntu 目录结构}\label{ubuntu-ux76eeux5f55ux7ed3ux6784}

{ % do not increment counter
\begin{longtable}[]{@{}ll@{}}
\toprule\noalign{}
目录 & 主要用途 \\
\midrule\noalign{}
\endhead
\bottomrule\noalign{}
\endlastfoot
\passthrough{\lstinline!/bin!} & 重要的二进制应用程序 \\
\passthrough{\lstinline!/boot!} & 启动配置文件和内核相关文件 \\
\passthrough{\lstinline!/dev!} & 设备文件，以文件形式表示设备和接口 \\
\passthrough{\lstinline!/etc!} & 系统配置文件、启动脚本和服务配置 \\
\passthrough{\lstinline!/home!} & 普通用户主目录 \\
\passthrough{\lstinline!/lib!} & 系统库文件 \\
\passthrough{\lstinline!/mnt!} & 临时挂载文件系统的位置 \\
\passthrough{\lstinline!/opt!} & 可选应用程序安装目录 \\
\passthrough{\lstinline!/usr/local!} & 用户自行安装的软件和运行环境 \\
\end{longtable}
}

\passthrough{\lstinline!/dev!}
中的文件不是普通文档。\passthrough{\lstinline!/dev/video0!}、\passthrough{\lstinline!/dev/snd!}
和 \passthrough{\lstinline!/dev/bus/usb!}
对应真实设备，访问它们会改变设备状态或与硬件交互。\passthrough{\lstinline!/etc!}
下的文件通常需要 root 权限，修改前应备份并确认目标主机。

\subsection{绝对路径、相对路径和工作目录}\label{ux7eddux5bf9ux8defux5f84ux76f8ux5bf9ux8defux5f84ux548cux5de5ux4f5cux76eeux5f55}

绝对路径从 \passthrough{\lstinline!/!} 开始，例如
\passthrough{\lstinline!/home/<user>/Documents/ascend310!}。相对路径从当前工作目录开始，例如
\passthrough{\lstinline!models/example.om!}。使用
\passthrough{\lstinline!..!} 返回上一层：

\begin{lstlisting}[language=bash]
cd ~/Documents
pwd
cd ..
pwd
cd -
pwd
\end{lstlisting}

\passthrough{\lstinline!cd!}
不带参数时返回用户主目录，\passthrough{\lstinline!cd -!}
返回上一个工作目录。进入服务、转换或同步命令之前，先执行
\passthrough{\lstinline!pwd!}，确认没有在错误目录中运行。

解析实际路径：

\begin{lstlisting}[language=bash]
realpath .
realpath models/example.om
readlink -f frontend/dist/index.html
\end{lstlisting}

\passthrough{\lstinline!realpath!} 和
\passthrough{\lstinline!readlink -f!}
适合在复制、删除、挂载和部署前确认符号链接真实指向。路径中包含空格时，变量和参数应使用双引号：

\begin{lstlisting}[language=bash]
target_dir="$HOME/Documents/board assets"
mkdir -p "$target_dir"
cp -a config.example.yaml "$target_dir/"
\end{lstlisting}

\subsection{创建、查看和复制}\label{ux521bux5efaux67e5ux770bux548cux590dux5236}

\begin{lstlisting}[language=bash]
mkdir -p reports/board/$(date --iso-8601=date)
touch notes.txt
ls -la
ls -lh models
cp -a config.example.yaml config.yaml
cp -r templates templates-backup
mv old-name.txt new-name.txt
\end{lstlisting}

\passthrough{\lstinline!mkdir -p!}
创建多级目录，\passthrough{\lstinline!cp -a!}
尽量保留文件属性，\passthrough{\lstinline!cp -r!}
复制目录，\passthrough{\lstinline!mv!}
同时用于移动和重命名。覆盖已有文件之前，先用
\passthrough{\lstinline!ls!}、\passthrough{\lstinline!test -e!} 或
\passthrough{\lstinline!stat!} 检查目标：

\begin{lstlisting}[language=bash]
test -e config.yaml && echo 'target exists'
stat -c '%A %U:%G %s %n' config.yaml
file config.yaml
\end{lstlisting}

\passthrough{\lstinline!stat!}
显示权限、所有者、字节数和名称，\passthrough{\lstinline!file!}
只根据文件内容给出格式线索，不能证明模型可以由 ACL 加载。

\subsection{查看文本和日志}\label{ux67e5ux770bux6587ux672cux548cux65e5ux5fd7}

\begin{lstlisting}[language=bash]
head -n 40 service.log
tail -n 80 service.log
tail -f service.log
wc -l service.log
sed -n '1,80p' service.log
awk '{print NR ":" $0}' service.log | tail -n 20
grep -n 'ERROR\|WARN' service.log
\end{lstlisting}

\passthrough{\lstinline!tail -f!}
持续跟踪日志，适合前台排障；批量检查应使用
\passthrough{\lstinline!head!}、\passthrough{\lstinline!tail!}、\passthrough{\lstinline!wc!}、\passthrough{\lstinline!sed!}
和 \passthrough{\lstinline!grep!}，并把输出保存到报告。仓库内搜索可使用
\passthrough{\lstinline!rg!}：

\begin{lstlisting}[language=bash]
rg -n 'atc|acl|sha256|127\.0\.0\.1|0\.0\.0\.0' README.md docs scripts
\end{lstlisting}

\subsection{查找、容量和文件类型}\label{ux67e5ux627eux5bb9ux91cfux548cux6587ux4ef6ux7c7bux578b}

\begin{lstlisting}[language=bash]
find . -maxdepth 2 -type f -print | sort
du -sh models data reports
df -h .
lsblk
\end{lstlisting}

\passthrough{\lstinline!find!} 先用 \passthrough{\lstinline!-maxdepth!}
限制范围，避免扫描整个家目录或挂载盘；\passthrough{\lstinline!du!}
查看目录占用，\passthrough{\lstinline!df!}
查看文件系统剩余空间，两者含义不同。\passthrough{\lstinline!lsblk!}
只用于查看块设备树，不能未经确认就挂载或格式化设备。

\subsection{归档和校验}\label{ux5f52ux6863ux548cux6821ux9a8c}

\begin{lstlisting}[language=bash]
run_date="$(date --iso-8601=date)"
tar -czf "reports/board/smoke-$run_date.tar.gz" reports/board/smoke.log
tar -tzf "reports/board/smoke-$run_date.tar.gz" | head
unzip -l artifact.zip
sha256sum models/example.om
\end{lstlisting}

压缩前先确认内容。不要把照片、生物特征模板、数据库、模型二进制、密钥和
token 放入公开归档。\passthrough{\lstinline!sha256sum!}
只能证明字节一致，不能证明模型正确、可加载或数值正确。

\section{4. 用户、组、权限、sudo
与安全删除}\label{ux7528ux6237ux7ec4ux6743ux9650sudo-ux4e0eux5b89ux5168ux5220ux9664}

\subsection{Linux 用户和组}\label{linux-ux7528ux6237ux548cux7ec4}

Ubuntu 安装时会创建一个普通用户，同时保留具有更高权限的 root
用户。普通用户在自己的主目录中通常不需要
\passthrough{\lstinline!sudo!}，访问
\passthrough{\lstinline!/etc!}、安装软件包或修改系统服务时才需要提权。

每个文件都有三组权限：

{ % do not increment counter
\begin{longtable}[]{@{}ll@{}}
\toprule\noalign{}
权限组 & 适用范围 \\
\midrule\noalign{}
\endhead
\bottomrule\noalign{}
\endlastfoot
user & 文件所有者 \\
group & 文件所属组 \\
other & 其他用户 \\
\end{longtable}
}

每组权限包含读 \passthrough{\lstinline!r!}、写
\passthrough{\lstinline!w!}、执行 \passthrough{\lstinline!x!}：

\begin{lstlisting}[language=bash]
ls -l script.sh
stat -c '%A %a %U:%G %n' script.sh
id
groups
\end{lstlisting}

\passthrough{\lstinline!ls -l!}
的权限字符串分为三段。\passthrough{\lstinline!rwxr-xr-x!}
表示所有者可读写执行，组和其他用户可读可执行。\passthrough{\lstinline!stat!}
的 \passthrough{\lstinline!\%a!}
给出八进制权限，\passthrough{\lstinline!\%U:\%G!} 给出所有者和组。

\subsection{修改权限}\label{ux4feeux6539ux6743ux9650}

给脚本增加所有者执行权限：

\begin{lstlisting}[language=bash]
chmod u+x scripts/run_service.sh
bash -n scripts/run_service.sh
./scripts/run_service.sh --help
\end{lstlisting}

\passthrough{\lstinline!chmod u+x!}
只给文件所有者增加执行权限。若脚本需要其他用户执行，应先确认运行用户和目录权限，再选择精确的组权限或
ACL。不要对整个家目录、模型目录或系统目录执行
\passthrough{\lstinline!chmod -R 777!}。

修改所有者和组需要管理员权限：

\begin{lstlisting}[language=bash]
sudo chown "$USER":"$USER" path/to/file
sudo chgrp developers path/to/directory
\end{lstlisting}

\passthrough{\lstinline!chown -R!} 和 \passthrough{\lstinline!chgrp -R!}
会递归改变整棵树，执行前必须用 \passthrough{\lstinline!find!}
预览目标，并确认没有系统文件和共享资产混在其中。

\subsection{sudo}\label{sudo}

\passthrough{\lstinline!sudo!} 表示以超级用户权限执行单条命令：

\begin{lstlisting}[language=bash]
sudo apt-get update
sudo systemctl status ssh --no-pager
sudo -i
\end{lstlisting}

\passthrough{\lstinline!sudo -i!} 进入 root
环境，\passthrough{\lstinline!sudo su!} 也可能加载 root
环境或保留部分当前环境。两者都不是日常编辑项目文件的默认方式。退出 root
状态使用：

\begin{lstlisting}[language=bash]
exit
\end{lstlisting}

频繁使用系统管理命令时，应通过明确的 sudo
策略和组权限管理，而不是把普通用户长期改成 root。

\subsection{设备组权限}\label{ux8bbeux5907ux7ec4ux6743ux9650}

摄像头和音频设备常通过
\passthrough{\lstinline!video!}、\passthrough{\lstinline!audio!}
等组授权：

\begin{lstlisting}[language=bash]
groups
ls -l /dev/video* /dev/snd
sudo usermod -aG video "$USER"
\end{lstlisting}

执行 \passthrough{\lstinline!usermod!}
后需要重新登录，组变更才会进入新会话。不要用
\passthrough{\lstinline!sudo python!}
启动整个服务来绕开设备权限，因为这会改变服务的数据目录所有者、环境变量和日志归属。

\subsection{安全删除}\label{ux5b89ux5168ux5220ux9664}

\passthrough{\lstinline!rm!}
不提供回收站。删除测试数据时，先把目标解析为绝对路径并验证：

\begin{lstlisting}[language=bash]
target_dir="$(realpath -- "$PALMPRINT_ROOT/data/captures")"
case "$target_dir" in
  "$PALMPRINT_ROOT"/data/captures) ;;
  *) printf 'refuse to remove: %s\n' "$target_dir" >&2; exit 1 ;;
esac
find "$target_dir" -maxdepth 1 -type f -print
rm -f -- "$target_dir"/*.jpg
\end{lstlisting}

先运行 \passthrough{\lstinline!find!}
检查文件列表，再由操作者确认后执行删除。\passthrough{\lstinline!rm -rf!}
只应在脚本创建的隔离临时目录中使用，并且变量必须非空、路径必须已经解析。不要把变量为空、路径未验证的命令变成
\passthrough{\lstinline!rm -rf "$HOME"!}、\passthrough{\lstinline!rm -rf .!}
或对整个部署目录的递归删除。

\section{5. nano 与 vim
编辑器}\label{nano-ux4e0e-vim-ux7f16ux8f91ux5668}

\subsection{nano}\label{nano}

打开文件：

\begin{lstlisting}[language=bash]
nano config.yaml
sudo nano /etc/exports
\end{lstlisting}

nano 进入后可以直接编辑文本。保存时按
\passthrough{\lstinline!Ctrl+O!}，屏幕会显示文件名；不修改文件名就按回车。退出按
\passthrough{\lstinline!Ctrl+X!}。如果文件有未保存修改，nano
会询问是否保存，按屏幕提示选择。

编辑系统文件前先备份：

\begin{lstlisting}[language=bash]
sudo cp /etc/apt/sources.list /etc/apt/sources.list.bak
sudo nano /etc/apt/sources.list
\end{lstlisting}

\subsection{vim}\label{vim}

打开文件：

\begin{lstlisting}[language=bash]
vim config.yaml
sudo vim ~/.config/autostart/myprogram.desktop
\end{lstlisting}

vim 初始处于普通模式。按 \passthrough{\lstinline!i!} 或
\passthrough{\lstinline!a!} 进入插入模式，按
\passthrough{\lstinline!Esc!} 回到普通模式。常用移动和编辑键：

{ % do not increment counter
\begin{longtable}[]{@{}ll@{}}
\toprule\noalign{}
按键 & 作用 \\
\midrule\noalign{}
\endhead
\bottomrule\noalign{}
\endlastfoot
\passthrough{\lstinline!k!} / \passthrough{\lstinline!j!} & 光标上移 /
下移 \\
\passthrough{\lstinline!h!} / \passthrough{\lstinline!l!} & 光标左移 /
右移 \\
\passthrough{\lstinline!i!} / \passthrough{\lstinline!I!} &
当前光标处插入 / 行首插入 \\
\passthrough{\lstinline!a!} / \passthrough{\lstinline!A!} & 光标后插入 /
行末插入 \\
\passthrough{\lstinline!o!} / \passthrough{\lstinline!O!} &
当前行后新建一行 / 当前行前新建一行 \\
\passthrough{\lstinline!x!} / \passthrough{\lstinline!X!} & 删除当前字符
/ 删除前一个字符 \\
\passthrough{\lstinline!dd!} & 剪切整行，也可作为删除 \\
\passthrough{\lstinline!dw!} & 删除一个单词 \\
\passthrough{\lstinline!d\^!} & 删除到行首 \\
\passthrough{\lstinline!dG!} / \passthrough{\lstinline!d1G!} &
删除到文档末尾 / 删除到文档开头 \\
\passthrough{\lstinline!gg!} / \passthrough{\lstinline!Shift+g!} &
移到第一行 / 最后一行 \\
\passthrough{\lstinline!yy!} / \passthrough{\lstinline!p!} /
\passthrough{\lstinline!P!} & 复制整行 / 粘贴到光标后 / 粘贴到光标前 \\
\end{longtable}
}

保存和退出前先按 \passthrough{\lstinline!Esc!}
回到普通模式，再输入冒号命令：

{ % do not increment counter
\begin{longtable}[]{@{}ll@{}}
\toprule\noalign{}
命令 & 作用 \\
\midrule\noalign{}
\endhead
\bottomrule\noalign{}
\endlastfoot
\passthrough{\lstinline!:q!} & 退出 \\
\passthrough{\lstinline":q!"} & 强制退出，不保存 \\
\passthrough{\lstinline!:w!} & 保存 \\
\passthrough{\lstinline!:wq!} & 保存并退出 \\
\passthrough{\lstinline":wq!"} & 强制保存并退出 \\
\passthrough{\lstinline!:w <文件路径>!} & 另存为 \\
\end{longtable}
}

编辑器中的 \passthrough{\lstinline":q!"}
会丢失未保存内容，只在确认不需要修改时使用。需要编辑 root 文件时使用
\passthrough{\lstinline!sudo nano!} 或
\passthrough{\lstinline!sudo vim!}，不要把来源不明的
\passthrough{\lstinline!tee!} 保存技巧直接粘贴到终端。保存后使用
\passthrough{\lstinline!head!}、\passthrough{\lstinline!grep!} 或
\passthrough{\lstinline!git diff!} 检查文件内容。

\section{6. APT、换源、Conda、pip 与 Python
路径}\label{aptux6362ux6e90condapip-ux4e0e-python-ux8defux5f84}

\subsection{APT 软件包管理}\label{apt-ux8f6fux4ef6ux5305ux7ba1ux7406}

APT 操作会改变系统状态。执行前记录 Ubuntu 版本和计划安装的包：

\begin{lstlisting}[language=bash]
cat /etc/os-release
sudo apt-get update
sudo apt-get install -y package-name
sudo apt-get install -f
\end{lstlisting}

\passthrough{\lstinline!apt-get update!}
获取软件包列表，\passthrough{\lstinline!install!}
安装软件包，\passthrough{\lstinline!install -f!}
修复依赖。安装完成后记录命令、版本和输出：

\begin{lstlisting}[language=bash]
dpkg -s python3-dev python3-pip 2>/dev/null | grep -E '^(Package|Status|Version):'
apt-cache policy package-name
\end{lstlisting}

不要为了通过文档示例而重复安装或升级 CANN 相关包。板端已经验证过的
CANN、驱动和 Python 环境应保持原样，变更前先建立回滚方案。

\subsection{换源}\label{ux6362ux6e90}

Ubuntu 的软件源配置通常位于
\passthrough{\lstinline!/etc/apt/sources.list!}，部分新版本还会使用
\passthrough{\lstinline!/etc/apt/sources.list.d/!} 下的
\passthrough{\lstinline!.sources!} 文件。先备份，再根据 Ubuntu
版本替换为合适的国内镜像：

\begin{lstlisting}[language=bash]
sudo cp -a /etc/apt/sources.list /etc/apt/sources.list.bak
sudo nano /etc/apt/sources.list
sudo apt-get update
\end{lstlisting}

换源时确认发行版代号与镜像仓库匹配。不要混用不同 Ubuntu
版本的仓库，不要把第三方 PPA
当作系统基础源。更新失败时先恢复备份并检查错误，不要把错误源继续保留。

\subsection{Conda 环境}\label{conda-ux73afux5883}

先加载 Conda shell 函数，再激活环境：

\begin{lstlisting}[language=bash]
source /usr/local/miniconda3/etc/profile.d/conda.sh
conda activate base
python --version
python -c 'import sys; print(sys.executable)'
\end{lstlisting}

若脚本开启了 \passthrough{\lstinline!set -u!}，而 Conda 或 CANN
环境脚本读取未定义变量，可以使用第 2 节的
\passthrough{\lstinline!load\_cann!} 模式，只在加载厂商脚本时临时关闭
\passthrough{\lstinline!-u!}，加载完成后恢复。

创建独立环境：

\begin{lstlisting}[language=bash]
conda create -n case9-acl-om python=3.9
conda activate case9-acl-om
python -c 'import sys; print(sys.executable)'
\end{lstlisting}

环境名称和 Python
版本必须来自项目的版本要求。不要用其他版本的同名环境替换已经验证的板端环境。

\subsection{pip 与 Python
解释器}\label{pip-ux4e0e-python-ux89e3ux91caux5668}

优先使用 \passthrough{\lstinline!python -m pip!}，让 pip
对应当前解释器：

\begin{lstlisting}[language=bash]
python -m pip --version
python -m pip install -r requirements.txt
python -m pip install -r requirements-board.txt
python -m pip list
\end{lstlisting}

不要使用 \passthrough{\lstinline!sudo pip install!} 或
\passthrough{\lstinline!sudo python!}。前者会把包安装到系统目录并绕过
Conda，后者会改变脚本的解释器和文件所有权。安装失败时先记录
\passthrough{\lstinline!python -m pip --version!}、\passthrough{\lstinline!which python!}、\passthrough{\lstinline!sys.executable!}
和完整错误，再决定是否变更环境。

\subsection{Python
路径和动态库路径}\label{python-ux8defux5f84ux548cux52a8ux6001ux5e93ux8defux5f84}

\begin{lstlisting}[language=bash]
export PYTHONNOUSERSITE=1
export PYTHONPATH="$PWD:${PYTHONPATH:-}"
export LD_LIBRARY_PATH="/usr/local/Ascend/ascend-toolkit/latest/lib64:${LD_LIBRARY_PATH:-}"
python -c 'import sys; print("\n".join(sys.path))'
python -c 'import os; print(os.environ.get("LD_LIBRARY_PATH", ""))'
\end{lstlisting}

\passthrough{\lstinline!PYTHONNOUSERSITE=1!}
用于避免用户目录中的包遮蔽板端环境，但可能让实际启动环境中存在的用户站点包不可见。判断导入问题时，应比较隔离环境和真实启动环境的
\passthrough{\lstinline!sys.executable!} 与
\passthrough{\lstinline!sys.path!}。\passthrough{\lstinline!PYTHONPATH!}
和 \passthrough{\lstinline!LD\_LIBRARY\_PATH!}
只加入已经确认的项目目录或 CANN
目录。持久化修改应经过部署评审，不能临时写入
\passthrough{\lstinline!/etc/profile!} 后让所有用户和所有服务继承。

\section{7. 网络、静态
IP、WiFi、热点与同网段检查}\label{ux7f51ux7edcux9759ux6001-ipwifiux70edux70b9ux4e0eux540cux7f51ux6bb5ux68c0ux67e5}

\subsection{网络状态}\label{ux7f51ux7edcux72b6ux6001}

\begin{lstlisting}[language=bash]
ip addr
ip route
nmcli connection show
nmcli -t -f NAME connection show
ping -c 4 <target-ip>
\end{lstlisting}

\passthrough{\lstinline!ip addr!}
查看接口地址，\passthrough{\lstinline!ip route!}
查看默认路由和网段，\passthrough{\lstinline!nmcli connection show!} 查看
NetworkManager 连接配置。\passthrough{\lstinline!ping!}
成功只说明网络层可达，不代表 SSH、VNC、NFS 或 NPU 服务已经可用。

\subsection{一般静态 IP
概念}\label{ux4e00ux822cux9759ux6001-ip-ux6982ux5ff5}

Ubuntu 图形桌面可以通过 NetworkManager 配置静态 IP。打开连接编辑器：

\begin{lstlisting}[language=bash]
nm-connection-editor
\end{lstlisting}

在对应连接中进入 \passthrough{\lstinline!IPv4 Settings!}，将方法设为
\passthrough{\lstinline!Manual!}，再填写地址、子网掩码、网关和
DNS。静态地址留在目标设备所在网段内，且不能与其他设备冲突。修改后重新启用该连接，再验证：

\begin{lstlisting}[language=bash]
ip -4 addr show
ip route
ping -c 4 <同网段设备IP>
\end{lstlisting}

静态 IP 不等于端口已经开放。SSH 服务、VNC 共享开关、NFS
导出和防火墙都需要单独检查。不要在不确定目标网段和网关的情况下直接覆盖原配置，先在连接编辑器中记录原值，或保存一份配置备份。

\subsection{同网段判断}\label{ux540cux7f51ux6bb5ux5224ux65ad}

IPv4 地址由网络地址和主机地址组成，子网掩码决定两部分的分界。以
\passthrough{\lstinline!255.255.255.0!} 为例，前 24 位是网络地址，后 8
位是主机地址。\passthrough{\lstinline!192.168.1.10!} 和
\passthrough{\lstinline!192.168.1.20!}
与同一子网掩码进行按位与运算后，网络地址都是
\passthrough{\lstinline!192.168.1.0!}，因此属于同一网段。\passthrough{\lstinline!192.168.0.100!}
与 \passthrough{\lstinline!192.168.1.100!} 在同一掩码下不属于同一网段。

排障时检查客户端和板端是否拿到同一网段的地址、是否连接到同一个热点或有线网络、网关是否一致。若客户端和板端存在多个网卡，应使用
\passthrough{\lstinline!ip route get <目标IP>!}
确认实际出口接口，不要只看某一个网卡的地址。

\subsection{创建 WiFi 热点}\label{ux521bux5efa-wifi-ux70edux70b9}

图形界面通常从右上角网络图标进入
\passthrough{\lstinline!Edit Connections!}，点击加号，选择
\passthrough{\lstinline!WiFi!}，再点击
\passthrough{\lstinline!Create!}。命令行方式可以打开连接编辑器：

\begin{lstlisting}[language=bash]
nm-connection-editor
\end{lstlisting}

在 \passthrough{\lstinline!Wi-Fi!} 选项卡中设置：

\begin{itemize}
\tightlist
\item
  \passthrough{\lstinline!SSID!}：热点名称，和连接配置名称
  \passthrough{\lstinline!connection name!} 不是一回事。
\item
  \passthrough{\lstinline!Mode!}：选择
  \passthrough{\lstinline!Hotspot!}。
\item
  \passthrough{\lstinline!Band!}：可选自动、2.4 GHz 或 5
  GHz。频段拥挤时固定频段可能更稳定。
\item
  \passthrough{\lstinline!Device!}：多网卡主机需要指定实际使用的无线网卡。
\item
  \passthrough{\lstinline!Wi-Fi Security!}：选择
  \passthrough{\lstinline!WPA\&WPA2 Personal!} 并设置密码。
\item
  \passthrough{\lstinline!IPv4 Settings!}：按需填写静态地址、掩码和网关。
\item
  \passthrough{\lstinline!IPv6!}：不需要时可选择忽略。
\end{itemize}

一个无线网卡通常不能同时作为热点又连接外部
WiFi。需要同时提供热点和连接上游网络时，应使用两个独立网卡，并分别确认设备名和驱动状态。

\subsection{热点与 WiFi
模式切换}\label{ux70edux70b9ux4e0e-wifi-ux6a21ux5f0fux5207ux6362}

图形界面中先断开当前热点或
WiFi，再选择要连接的网络。重新开启热点时，从网络设置菜单选择热点配置并启用。切换后检查：

\begin{lstlisting}[language=bash]
ip addr
iw dev
nmcli device status
nmcli connection show --active
\end{lstlisting}

\passthrough{\lstinline!iw dev!}
是否可用取决于系统是否安装相关工具；即使命令不存在，也应记录
\passthrough{\lstinline!nmcli device status!}
和界面状态，而不是直接推断网卡损坏。切换过程中，已经连接的 SSH
会话可能短暂断开。使用无屏幕板端时，先准备好重新发现的 IP
的方法，再修改网络模式。

\subsection{自动连接优先级}\label{ux81eaux52a8ux8fdeux63a5ux4f18ux5148ux7ea7}

热点和 WiFi 配置可以在 \passthrough{\lstinline!General!} 选项卡中设置
\passthrough{\lstinline!Connect automatically with priority!}。不同版本的教程对数值方向描述不一致：有的资料写``数值越大优先级越高''，有的资料写``数值越小优先级越高''。不要只照抄数字，应查看当前
NetworkManager 的实际配置并做一次重启验证：

\begin{lstlisting}[language=bash]
nmcli -f connection.id,connection.autoconnect,connection.autoconnect-priority connection show
\end{lstlisting}

记录哪一个连接在重启后实际被选中，再据此调整优先级。修改网络配置后，若板端没有屏幕，先确认仍能通过已有热点或串口获取
IP，防止失去远程入口。

\subsection{无屏幕联网的基本流程}\label{ux65e0ux5c4fux5e55ux8054ux7f51ux7684ux57faux672cux6d41ux7a0b}

无屏幕板端的基本顺序是：默认热点可用，客户端连接热点，通过 SSH
登录，扫描可用 WiFi，选择目标 SSID 并输入密码，获取新的
IP，然后让客户端连接到同一
WiFi。切换时可能短暂断开当前热点，需要准备重新连接或通过串口查看 IP。

\begin{lstlisting}[language=bash]
ssh -Y <user>@<board-ip>
nmcli device wifi list
nmcli connection show
ip addr
\end{lstlisting}

具体热点名称、默认 IP
和脚本路径必须来自当前板端的实际配置。不要把其他产品的示例热点名或脚本路径当作
310B 板的默认值。成功连接 WiFi 后，IP 可能改变；应记录 DHCP
分配的地址或为板端设置不冲突的静态地址。

\section{8. SSH
远程命令行控制}\label{ssh-ux8fdcux7a0bux547dux4ee4ux884cux63a7ux5236}

\subsection{客户端与服务端}\label{ux5ba2ux6237ux7aefux4e0eux670dux52a1ux7aef}

SSH 使用客户端和服务端模型。310B 板端作为服务端运行
\passthrough{\lstinline!sshd!}，Ubuntu
主机作为客户端发起连接。客户端连接后看到的 shell 就是板端
shell，因此文件路径、权限、设备和运行时环境都以板端为准。

\subsection{安装和启动服务端}\label{ux5b89ux88c5ux548cux542fux52a8ux670dux52a1ux7aef}

在板端确认 SSH 状态：

\begin{lstlisting}[language=bash]
sudo service ssh status
\end{lstlisting}

未安装时执行：

\begin{lstlisting}[language=bash]
sudo apt update
sudo apt install openssh-server
sudo service ssh start
sudo systemctl enable ssh
\end{lstlisting}

再次检查：

\begin{lstlisting}[language=bash]
systemctl status ssh --no-pager
ss -ltnp | grep ':22'
\end{lstlisting}

Ubuntu 主机通常自带 SSH 客户端。需要单独安装时执行：

\begin{lstlisting}[language=bash]
sudo apt-get install openssh-client
ssh -V
\end{lstlisting}

\subsection{网络连通性检查}\label{ux7f51ux7edcux8fdeux901aux6027ux68c0ux67e5}

先连接板端 WiFi 或有线网络，再检查同网段和可达性：

\begin{lstlisting}[language=bash]
ip addr
ip route
ping -c 4 192.168.0.100
\end{lstlisting}

若无法 ping
通，依次检查：两台设备是否在同一网段、客户端是否连接了错误的网络、板端
IP
是否已经变化、路由出口是否正确、无线热点是否仍处于开启状态。\passthrough{\lstinline!ping!}
不通时不要直接反复执行
\passthrough{\lstinline!ssh!}，先把网络层问题定位清楚。

\subsection{登录}\label{ux767bux5f55}

源文档中的典型命令是：

\begin{lstlisting}[language=bash]
ssh -Y wheeltec@192.168.0.100
\end{lstlisting}

其中 \passthrough{\lstinline!-Y!} 启用受信任的 X11
转发，\passthrough{\lstinline!wheeltec!}
是服务端用户名，\passthrough{\lstinline!192.168.0.100!}
是服务端地址。实际使用时替换为已经确认的值，不要复制示例密码。仅需要命令行时可以使用普通
SSH：

\begin{lstlisting}[language=bash]
ssh <user>@<board-ip>
ssh <user>@<board-ip> 'hostname; uname -a; npu-smi info'
\end{lstlisting}

登录后提示符中的用户名和主机名会变化。若用户名仍是本地用户，说明没有进入远端
shell。需要执行单条板端命令时，把命令放在引号中，并确认引号内的变量是在远端还是本地展开。

\subsection{known\_hosts 冲突}\label{known_hosts-ux51b2ux7a81}

第一次连接某台设备或设备身份发生变化时，可能出现主机密钥冲突。源文档给出的处理命令是：

\begin{lstlisting}[language=bash]
ssh-keygen -f "/home/wheeltec-client/.ssh/known_hosts" -R "192.168.0.100"
ssh -Y wheeltec@192.168.0.100
\end{lstlisting}

第一条命令只删除指定主机在 \passthrough{\lstinline!known\_hosts!}
中的记录，第二条重新连接。执行删除前先确认 IP
确实属于目标板端，并核对新的主机密钥指纹。不要为了省事直接删除整个
\passthrough{\lstinline!known\_hosts!}，也不要在未确认设备身份时盲目确认新密钥。

\subsection{同一网段排障}\label{ux540cux4e00ux7f51ux6bb5ux6392ux969c}

如果 SSH 连接失败，检查客户端和板端的 IPv4 地址与子网掩码：

\begin{lstlisting}[language=bash]
ip -4 addr show
ip route
ping -c 4 <board-ip>
\end{lstlisting}

静态 IP
必须位于同一网段，且不能与板端地址或其他设备冲突。若修改的是客户端地址，保留原配置；若修改的是板端地址，先确认自己有第二条访问通道。路由器、热点、交换机和防火墙都可能影响可达性。

\subsection{SSH
安全注意事项}\label{ssh-ux5b89ux5168ux6ce8ux610fux4e8bux9879}

\begin{itemize}
\tightlist
\item
  使用强密码或 SSH 密钥，优先使用密钥登录。
\item
  只允许可信网络访问 22 端口，不要把 SSH 直接暴露到公网。
\item
  需要远程管理时，先确认服务端主机密钥指纹，再建立连接。
\item
  \passthrough{\lstinline!ssh -Y!} 会转发 X11
  连接，信任边界较宽，只在受信网络中使用。
\item
  多人共用板端时，为每个用户分配账号，不要共享 root 密码。
\item
  修改防火墙或 SSH
  配置后，保留一个已登录会话，再从新会话验证，避免把自己锁在系统外。
\end{itemize}

\section{9. VNC
远程桌面控制}\label{vnc-ux8fdcux7a0bux684cux9762ux63a7ux5236}

\subsection{客户端与服务端}\label{ux5ba2ux6237ux7aefux4e0eux670dux52a1ux7aef-1}

VNC 也采用客户端和服务端模型。板端提供桌面服务，Ubuntu 主机或 Windows
主机运行 VNC
客户端。远程桌面适合查看图形窗口、图像预览和需要鼠标交互的程序；纯命令行任务优先使用
SSH。

\subsection{Ubuntu 客户端安装
Remmina}\label{ubuntu-ux5ba2ux6237ux7aefux5b89ux88c5-remmina}

源文档给出的 Remmina 安装命令是：

\begin{lstlisting}[language=bash]
sudo apt-add-repository ppa:remmina-ppa-team/remmina-next
sudo apt update
sudo apt install remmina remmina-plugin-rdp remmina-plugin-secret
\end{lstlisting}

PPA 是否适合当前 Ubuntu
版本需要先确认。若系统仓库已经提供合适版本，也可以只安装
\passthrough{\lstinline!remmina!}。安装完成后从应用菜单启动，选择 VNC
协议，填写板端 IP 地址，再输入 VNC 密码。源文档示例 IP 为
\passthrough{\lstinline!192.168.0.100!}，实际使用时必须替换。

\subsection{Windows 客户端使用
MobaXterm}\label{windows-ux5ba2ux6237ux7aefux4f7fux7528-mobaxterm}

在 Windows 主机上启动 MobaXterm，点击左上角
\passthrough{\lstinline!Session!}，选择
\passthrough{\lstinline!VNC!}，填写板端 IP
地址，确认客户端已经连接到板端所在网络，然后点击
\passthrough{\lstinline!OK!} 并输入 VNC
密码。不同版本的菜单名称可能略有变化，但客户端、目标
IP、密码和网络可达性这几项不变。

\subsection{分辨率和画面质量}\label{ux5206ux8fa8ux7387ux548cux753bux9762ux8d28ux91cf}

板端没有连接显示器时，VNC 分辨率可能较低。源文档给出的调整命令是：

\begin{lstlisting}[language=bash]
xrandr --fb 1024x768
\end{lstlisting}

\passthrough{\lstinline!xrandr --fb!} 设置当前 X
显示的帧缓冲尺寸。执行前先运行 \passthrough{\lstinline!xrandr!}
查看可用输出和模式，确认目标分辨率不会超出显示能力。若画面卡顿，可以在
Remmina
中降低画质、颜色深度或更新频率。图像质量和实时性需要权衡，不要把降低画质后的流畅度当作原始编码性能。

\subsection{共享开关}\label{ux5171ux4eabux5f00ux5173}

部分 Ubuntu 桌面镜像使用共享设置提供 VNC。修改网线、IP 或网络模式后，VNC
可能断开，需要检查桌面设置中的 \passthrough{\lstinline!Sharing!} 和
\passthrough{\lstinline!Screen Sharing!}
是否仍为活动状态。带有远程桌面协议的版本可能还需要启用兼容 VNC
协议的选项。共享开关关闭时，SSH 可能仍正常，但 VNC 无法连接。

\subsection{屏幕和显示要求}\label{ux5c4fux5e55ux548cux663eux793aux8981ux6c42}

部分板端镜像在没有连接显示器时不会生成可用的桌面帧缓冲。VNC
连接前可能需要连接真实屏幕或使用显示欺骗器。无屏幕场景应先确认 VNC
服务、共享开关、显示输出和分辨率，再判断是否为客户端问题。

\subsection{VNC 中打开终端}\label{vnc-ux4e2dux6253ux5f00ux7ec8ux7aef}

大部分 VNC 客户端无法让 \passthrough{\lstinline!Ctrl+Alt+T!}
正确传递到远端桌面。源文档建议使用鼠标右键菜单中的终端入口。打开终端后，先用
\passthrough{\lstinline!whoami!}、\passthrough{\lstinline!hostname!} 和
\passthrough{\lstinline!pwd!} 确认会话确实运行在板端。

\subsection{SSH 与 VNC 的对比}\label{ssh-ux4e0e-vnc-ux7684ux5bf9ux6bd4}

{ % do not increment counter
\begin{longtable}[]{@{}lll@{}}
\toprule\noalign{}
对比项 & SSH & VNC \\
\midrule\noalign{}
\endhead
\bottomrule\noalign{}
\endlastfoot
传输内容 & 终端文本和可选 X11 转发 & 桌面图像、鼠标和键盘事件 \\
带宽占用 & 通常较低 & 通常较高，受分辨率和画质影响 \\
适合任务 & 命令、日志、脚本、服务管理 & 图形程序、图像窗口、桌面交互 \\
无屏幕要求 & 不要求桌面帧缓冲 & 部分镜像要求显示器或显示欺骗器 \\
安全边界 & 可使用密钥和严格访问控制 & 需要保护 VNC 密码和端口 \\
断线影响 & 前台命令可能随会话结束 & 前台图形程序可能随会话结束 \\
\end{longtable}
}

VNC 服务不要直接暴露到不可信网络。需要通过 SSH
隧道访问时，可以把本地端口转发到板端 VNC 端口，但应确认 VNC
监听地址和端口，并使用唯一强密码。

\section{10. NFS、SCP 与 rsync}\label{nfsscp-ux4e0e-rsync}

\subsection{NFS 的角色}\label{nfs-ux7684ux89d2ux8272}

NFS 把服务端目录映射到客户端挂载点。典型的板端作为服务端，Ubuntu
主机作为客户端。挂载后，客户端编辑文件就像编辑本地目录，但底层文件仍属于远端文件系统，断网、服务端关机或网络切换会造成挂载失效。

\subsection{服务端安装和导出}\label{ux670dux52a1ux7aefux5b89ux88c5ux548cux5bfcux51fa}

板端安装服务端：

\begin{lstlisting}[language=bash]
sudo apt-get install nfs-kernel-server
\end{lstlisting}

编辑导出配置：

\begin{lstlisting}[language=bash]
sudo nano /etc/exports
\end{lstlisting}

文件内容示例：

\begin{lstlisting}
/srv/ascend310-share *(rw,sync,no_root_squash)
\end{lstlisting}

第一列是共享目录，第二列是允许访问的客户端范围。生产或多人环境中应把
\passthrough{\lstinline!*!} 收窄为具体网段或主机地址，并根据需要改为
\passthrough{\lstinline!ro!}。\passthrough{\lstinline!no\_root\_squash!}
会让客户端 root 保留 root
权限，安全风险较高；只有在隔离实验网络中经过评估后才使用。

源文档还出现过下面的旧示例：

\begin{lstlisting}
/home/wheeltec/wheeltec_ros2 *(rw,sync,no_root_squash)
\end{lstlisting}

该路径只是示例，不是 310B
板的默认路径。使用时必须替换为实际共享目录，并确认目录存在、内容范围和客户端权限。

启动服务：

\begin{lstlisting}[language=bash]
sudo /etc/init.d/nfs-kernel-server start
sudo /etc/init.d/nfs-kernel-server restart
systemctl status nfs-kernel-server --no-pager
exportfs -v
\end{lstlisting}

\passthrough{\lstinline!exportfs -v!}
查看当前导出，\passthrough{\lstinline!systemctl status!}
查看服务状态。修改 \passthrough{\lstinline!/etc/exports!}
后通常需要重新加载或重启服务。

\subsection{客户端安装和挂载}\label{ux5ba2ux6237ux7aefux5b89ux88c5ux548cux6302ux8f7d}

Ubuntu 客户端安装：

\begin{lstlisting}[language=bash]
sudo apt-get install nfs-common
\end{lstlisting}

创建挂载点：

\begin{lstlisting}[language=bash]
sudo mkdir -p /mnt/mount_nfs
\end{lstlisting}

先确认可达，再挂载：

\begin{lstlisting}[language=bash]
ping -c 4 <server-ip>
sudo mount -t nfs <server-ip>:/srv/ascend310-share /mnt/mount_nfs
mount | grep /mnt/mount_nfs
\end{lstlisting}

需要指定 NFS 版本或锁选项时，使用明确的参数：

\begin{lstlisting}[language=bash]
sudo mount -t nfs -o nolock <server-ip>:/srv/ascend310-share /mnt/mount_nfs
\end{lstlisting}

卸载：

\begin{lstlisting}[language=bash]
sudo umount -t nfs <server-ip>:/srv/ascend310-share /mnt/mount_nfs
\end{lstlisting}

若服务端曾经关机，挂载可能变成失效状态，重新挂载前先卸载。不要对正在写入的共享目录执行
\passthrough{\lstinline!umount -f!}，除非已经确认没有进程使用并且接受数据风险。

\subsection{NFS 权限和安全}\label{nfs-ux6743ux9650ux548cux5b89ux5168}

源文档中的旧配置使用了宽泛权限。下面的命令只作为说明，不建议直接照抄：

\begin{lstlisting}
sudo chmod -R 777 /path/to/share
sudo chown -R 777 nobody /path/to/share
\end{lstlisting}

\passthrough{\lstinline!chmod -R 777!}
会向所有用户开放读写执行；\passthrough{\lstinline!chown -R 777 nobody!}
不是通用安全的用户和组写法。更合理的做法是创建专用组、设置
\passthrough{\lstinline!u+rwX!}
或组权限、把客户端限制在同一网段，并在不需要写入时使用只读导出。模型、数据、模板和报告目录应按实际所有者设置权限，不要用一次宽泛递归命令处理整个家目录。

\subsection{SCP 定向传输}\label{scp-ux5b9aux5411ux4f20ux8f93}

SCP 使用 SSH 传输文件：

\begin{lstlisting}[language=bash]
board_user="replace-with-board-user"
board_ip="replace-with-board-ip"
ssh_target="$board_user@$board_ip"
scp demo/vtest.avi "$ssh_target:Documents/ascend310/demo/"
scp "$ssh_target:Documents/ascend310/reports/result.json" .
\end{lstlisting}

下载和上传前都先检查目标目录和磁盘空间：

\begin{lstlisting}[language=bash]
ssh "$ssh_target" 'df -h "$HOME"; test -d "$HOME/Documents/ascend310/demo"'
\end{lstlisting}

\subsection{rsync 同步}\label{rsync-ux540cux6b65}

rsync 适合目录同步和断点续传。先预览：

\begin{lstlisting}[language=bash]
rsync -av --dry-run samples/case8/ \
  "$ssh_target:Documents/ascend310/samples/case8/"
\end{lstlisting}

确认输出后再执行：

\begin{lstlisting}[language=bash]
rsync -av --protect-args \
  samples/case8/ \
  "$ssh_target:Documents/ascend310/samples/case8/"
\end{lstlisting}

默认不要使用
\passthrough{\lstinline!--delete!}。它会删除远端已有文件，可能误删模型、数据或报告。需要排除大资产时：

\begin{lstlisting}[language=bash]
rsync -av --protect-args \
  --exclude 'data/' --exclude 'models/' --exclude '*.om' \
  samples/case8/ \
  "$ssh_target:Documents/ascend310/samples/case8/"
\end{lstlisting}

同步后分别在本地和远端执行
\passthrough{\lstinline!wc -c!}、\passthrough{\lstinline!sha256sum!}
或清单检查。模型、OM、数据集和真实运行报告不应因为一次同步就进入 Git。

\section{11. Docker 基础与 ROS2
容器命令}\label{docker-ux57faux7840ux4e0e-ros2-ux5bb9ux5668ux547dux4ee4}

\subsection{Docker 基本概念}\label{docker-ux57faux672cux6982ux5ff5}

Docker
使用镜像和容器组织应用。镜像包含文件系统和运行环境，容器是镜像的一次隔离运行实例。容器与宿主系统共享内核，但可以使用独立的文件系统、网络和进程空间。不要把容器当作可以随意获得所有硬件权限的沙箱。

检查 Docker 是否安装和运行：

\begin{lstlisting}[language=bash]
docker --version
systemctl status docker --no-pager
sudo systemctl start docker
sudo systemctl enable docker
\end{lstlisting}

查看容器和镜像：

\begin{lstlisting}[language=bash]
docker ps
docker ps -a
docker images
\end{lstlisting}

\passthrough{\lstinline!docker ps!}
只显示运行中的容器，\passthrough{\lstinline!docker ps -a!}
显示包括已停止的容器，\passthrough{\lstinline!docker images!}
显示本地镜像。

\subsection{通用容器生命周期命令}\label{ux901aux7528ux5bb9ux5668ux751fux547dux5468ux671fux547dux4ee4}

\begin{lstlisting}[language=bash]
docker start <container-name>
docker exec -it <container-name> bash
exit
docker stop <container-name>
docker restart <container-name>
docker update --restart=always <container-name>
docker inspect <container-name>
\end{lstlisting}

\passthrough{\lstinline!docker exec -it!}
进入正在运行的容器，\passthrough{\lstinline!exit!} 退出容器 shell
返回宿主。\passthrough{\lstinline!docker update --restart=always!}
会让容器随 Docker
服务自动启动，设置前应确认应用具备断电恢复和安全停机策略。

\subsection{教程中的 ROS2
容器命令}\label{ux6559ux7a0bux4e2dux7684-ros2-ux5bb9ux5668ux547dux4ee4}

部分镜像提供了封装脚本，树莓派 5 的 ROS2 示例命令名称为：

\begin{lstlisting}[language=bash]
ros2_start
ros2
ros2_stop
ros2_restart
\end{lstlisting}

这些名称是镜像提供的帮助命令或别名，不是 Docker 自带子命令。执行
\passthrough{\lstinline!ros2\_start!}
启动容器，\passthrough{\lstinline!ros2!}
进入容器环境，\passthrough{\lstinline!ros2\_stop!}
停止容器，\passthrough{\lstinline!ros2\_restart!}
重启容器。使用前先检查它们是脚本、别名还是函数：

\begin{lstlisting}[language=bash]
type ros2_start
type ros2
type ros2_stop
type ros2_restart
\end{lstlisting}

进入容器后看到的用户提示符可能变化。容器状态异常时先查看
\passthrough{\lstinline!docker ps -a!}、\passthrough{\lstinline!docker logs <container-name>!}
和
\passthrough{\lstinline!docker inspect!}，不要直接删除容器或镜像。ROS2
节点、话题、launch 和工作空间编译属于 ROS2 内容，应放在附录 5 中讲解。

\subsection{源码目录、挂载和编译}\label{ux6e90ux7801ux76eeux5f55ux6302ux8f7dux548cux7f16ux8bd1}

容器可以把宿主目录挂载到容器路径。源文档提示容器内的源码可能已经映射到宿主目录，因此可以在宿主或容器中修改文件，但编译必须在安装了对应构建系统的容器中完成。使用前确认挂载关系：

\begin{lstlisting}[language=bash]
docker inspect -f '{{json .Mounts}}' <container-name>
docker exec -it <container-name> pwd
docker exec -it <container-name> ls -la
\end{lstlisting}

不要把宿主
\passthrough{\lstinline!/!}、\passthrough{\lstinline!/etc!}、Docker
socket 或整个家目录挂载到不可信容器。需要摄像头、串口、USB 或 NPU
设备时，逐项确认设备节点和权限，不要直接使用
\passthrough{\lstinline!--privileged!}。

\subsection{图形窗口和
DISPLAY}\label{ux56feux5f62ux7a97ux53e3ux548c-display}

容器中的图形程序需要显示服务器。源文档中的常见传递方式是：

\begin{lstlisting}[language=bash]
DISPLAY=${DISPLAY:-unix:0}
docker exec -it -e DISPLAY=$DISPLAY <container-id> bash
\end{lstlisting}

\passthrough{\lstinline!DISPLAY!}
是图形显示目标，\passthrough{\lstinline!-e!} 把它传给容器。使用 X11
转发时还要确认宿主端的 \passthrough{\lstinline!xhost!}
和认证设置。源教程使用：

\begin{lstlisting}
xhost +
\end{lstlisting}

该命令允许所有客户端连接显示服务器，范围过宽。受信多用户环境中应改为只允许当前本地用户，并在使用后恢复限制。SSH
会话的 \passthrough{\lstinline!DISPLAY!} 可能是
\passthrough{\lstinline!localhost:10.0!} 等转发地址，容器启动后新建的
SSH 会话不一定继承原有显示权限。

\section{12. 程序开机自启动：桌面项、启动应用和
systemd}\label{ux7a0bux5e8fux5f00ux673aux81eaux542fux52a8ux684cux9762ux9879ux542fux52a8ux5e94ux7528ux548c-systemd}

\subsection{桌面自启动项}\label{ux684cux9762ux81eaux542fux52a8ux9879}

桌面环境可以读取
\passthrough{\lstinline!\~/.config/autostart/*.desktop!} 自动启动程序：

\begin{lstlisting}[language=bash]
mkdir -p ~/.config/autostart
vim ~/.config/autostart/myprogram.desktop
\end{lstlisting}

文件内容：

\begin{lstlisting}
[Desktop Entry]
Encoding=UTF-8
Type=Application
Name=myprogram
Exec=/home/wheeltec/script.sh
Terminal=true
\end{lstlisting}

\passthrough{\lstinline!Exec!}
必须指向已经确认的脚本或程序，\passthrough{\lstinline!Terminal=true!}
表示需要打开终端。实际使用时把
\passthrough{\lstinline!/home/wheeltec/script.sh!}
替换为当前用户可执行的路径，并设置文件权限：

\begin{lstlisting}[language=bash]
chmod a+r ~/.config/autostart/myprogram.desktop
\end{lstlisting}

用户级桌面项只适用于图形登录用户。服务账号、无桌面会话和系统级程序应使用
systemd，而不是依赖桌面自启动。

\subsection{启动脚本}\label{ux542fux52a8ux811aux672c}

一个最小脚本应明确 shebang、工作目录和环境：

\begin{lstlisting}[language=bash]
#!/bin/bash
set -euo pipefail
cd /home/<user>/Documents/ascend310
exec /home/<user>/Documents/ascend310/scripts/run_service.sh
\end{lstlisting}

赋予所有者执行权限即可：

\begin{lstlisting}[language=bash]
chmod u+x script.sh
bash -n script.sh
./script.sh
\end{lstlisting}

源教程中的桌面脚本会调用终端模拟器并加载应用程序环境。若脚本只做桌面显示，\passthrough{\lstinline!Terminal=true!}
可以保留终端；若脚本管理后台服务，使用 systemd
并在服务单元中显式设置环境变量。

\subsection{Startup Applications}\label{startup-applications}

部分 Ubuntu 桌面提供 \passthrough{\lstinline!Startup Applications!}
图形工具。搜索 \passthrough{\lstinline!startup!} 打开该工具，点击
\passthrough{\lstinline!Add!}，填写名称和启动命令，保存后重启验证。它本质上仍是用户会话级自启动，不等价于系统服务。

\subsection{systemd 服务}\label{systemd-ux670dux52a1}

无桌面或需要自动重启的程序使用 systemd。创建单元文件：

\begin{lstlisting}[language=bash]
sudo vim /etc/systemd/system/ascend310-demo.service
\end{lstlisting}

示例内容：

\begin{lstlisting}
[Unit]
Description=Ascend 310B demo service
After=network-online.target
Wants=network-online.target

[Service]
Type=simple
User=<board-user>
WorkingDirectory=/home/<board-user>/Documents/ascend310
Environment=PYTHONNOUSERSITE=1
ExecStart=/home/<board-user>/Documents/ascend310/scripts/run_service.sh
Restart=on-failure
RestartSec=5

[Install]
WantedBy=multi-user.target
\end{lstlisting}

加载并启用：

\begin{lstlisting}[language=bash]
sudo systemctl daemon-reload
sudo systemctl enable --now ascend310-demo.service
systemctl status ascend310-demo.service --no-pager
journalctl -u ascend310-demo.service -b --no-pager | tail -n 80
\end{lstlisting}

systemd 不会自动继承交互式 shell 中临时 \passthrough{\lstinline!source!}
的环境。CANN、Conda、\passthrough{\lstinline!PYTHONPATH!} 和
\passthrough{\lstinline!LD\_LIBRARY\_PATH!}
应由包装脚本或服务单元显式设置。修改单元后重新执行
\passthrough{\lstinline!daemon-reload!}，停止或重启前先确认服务 PID
和影响范围。

\subsection{旧式 rc.local}\label{ux65e7ux5f0f-rc.local}

部分旧教程使用 \passthrough{\lstinline!/etc/rc.local!}：

\begin{lstlisting}[language=bash]
sudo nano /etc/rc.local
\end{lstlisting}

在文件末尾添加脚本路径，例如：

\begin{lstlisting}
/home/wheeltec/nfs.sh
\end{lstlisting}

然后重启：

\begin{lstlisting}[language=bash]
reboot
\end{lstlisting}

新版本 Ubuntu 默认不一定启用
\passthrough{\lstinline!rc.local!}，而且它不提供服务依赖、重启策略和日志单位。除非目标系统已有明确要求，优先使用
systemd。

\section{13. 昇腾 310B 的 CANN、PyACL、NPU、ATC 与
OM}\label{ux6607ux817e-310b-ux7684-cannpyaclnpuatc-ux4e0e-om}

\subsection{在同一 shell 加载 Conda 和
CANN}\label{ux5728ux540cux4e00-shell-ux52a0ux8f7d-conda-ux548c-cann}

\begin{lstlisting}[language=bash]
source /usr/local/miniconda3/etc/profile.d/conda.sh
conda activate base
source /usr/local/Ascend/ascend-toolkit/set_env.sh
python --version
python -c 'import sys; print(sys.executable)'
command -v atc
python -c 'import acl; print("PyACL import: ok")'
\end{lstlisting}

\passthrough{\lstinline!set\_env.sh!} 设置
ACL、运行时和工具链相关路径，\passthrough{\lstinline!import acl!}
验证当前解释器能否加载 PyACL。另一个终端中执行
\passthrough{\lstinline!source!} 不会影响当前服务进程。启动 ATC 或 NPU
服务时，Conda 和 CANN 必须在同一个 shell 或同一个包装脚本中加载。

\subsection{NPU 设备状态}\label{npu-ux8bbeux5907ux72b6ux6001}

\begin{lstlisting}[language=bash]
npu-smi info
\end{lstlisting}

\passthrough{\lstinline!npu-smi info!} 输出设备、芯片和健康状态。310B4 /
8T 与 310B1 / 20T 属于不同算力层级，结果不能混合排名。已知板端出现
\passthrough{\lstinline!Health: Alarm!}
时，应把它作为诊断背景保存，但不能仅凭告警判定
ATC、ACL、精度或性能失败；相反，应继续检查 \passthrough{\lstinline!acl!}
导入失败、设备不存在、ATC 非零退出、OM
缺失、推理不一致、段错误和资源泄漏等具体现象。

\subsection{ATC 转换}\label{atc-ux8f6cux6362}

ATC 把 ONNX 等前端模型转换为 OM。下面是最小示例：

\begin{lstlisting}[language=bash]
atc \
  --model=models/resnet18_scene.onnx \
  --framework=5 \
  --output=models/resnet18_scene \
  --soc_version=Ascend310B4 \
  --input_format=NCHW \
  --input_shape=input:1,3,224,224
\end{lstlisting}

\passthrough{\lstinline!--framework=5!} 表示
ONNX，\passthrough{\lstinline!--soc\_version=Ascend310B4!}
必须与目标板匹配，输入布局和 \passthrough{\lstinline!--input\_shape!}
必须来自该模型的输入合同。不要从另一个模型复制形状、算子参数或
\passthrough{\lstinline!--soc\_version!}。转换前先阅读项目脚本的
\passthrough{\lstinline!--help!}：

\begin{lstlisting}[language=bash]
python scripts/convert_onnx_to_om.py --help
SOC_VERSION=Ascend310B4 bash scripts/atc_convert.sh
\end{lstlisting}

脚本可能写入日志、manifest
和中间目录。执行前确认输出路径，执行后保存完整 ATC 命令、CANN
版本、模型输入合同、退出码和日志。

\subsection{OM 工件检查}\label{om-ux5de5ux4ef6ux68c0ux67e5}

转换完成后检查文件存在、字节数和摘要：

\begin{lstlisting}[language=bash]
stat -c '%s %n' models/*.om
sha256sum models/example.om
\end{lstlisting}

若复现包提供了
\passthrough{\lstinline!SHA256SUMS.txt!}，只能在生成该清单的根目录执行：

\begin{lstlisting}[language=bash]
bundle_root=/path/to/repro-bundle
cd "$bundle_root"
sha256sum -c SHA256SUMS.txt
\end{lstlisting}

摘要一致只证明文件没有发生变化。OM 是否能够由 PyACL
加载、输入输出是否与合同一致、数值是否可接受，需要分别执行 ACL
烟测和数值比较。不同 CANN 版本、不同 SoC、不同精度和不同模型 ID 的 OM
不能混用。

\subsection{PyACL 与 OM 推理}\label{pyacl-ux4e0e-om-ux63a8ux7406}

PyACL 是 ACL 的 Python 接口。服务启动前执行最小导入检查：

\begin{lstlisting}[language=bash]
python -c 'import acl; print("PyACL import: ok")'
\end{lstlisting}

推荐使用项目包装脚本，因为包装脚本可以在启动前检查
\passthrough{\lstinline!import acl!}：

\begin{lstlisting}[language=bash]
bash scripts/run_palmprint_service.sh --host 127.0.0.1 --port 7860
\end{lstlisting}

若 \passthrough{\lstinline!import acl!} 失败，记录 Python 路径、CANN
路径和完整错误，修复环境后重新启动服务。不要添加 CPU
推理回退，也不要用随机输出或模拟结果让健康检查通过。服务健康接口可以证明
HTTP 进程存活，不能证明 NPU 推理链已加载。

\subsection{环境变量边界}\label{ux73afux5883ux53d8ux91cfux8fb9ux754c}

\begin{lstlisting}[language=bash]
export PYTHONNOUSERSITE=1
export PYTHONPATH="$PWD:${PYTHONPATH:-}"
export LD_LIBRARY_PATH="/usr/local/Ascend/ascend-toolkit/latest/lib64:${LD_LIBRARY_PATH:-}"
\end{lstlisting}

这些变量只针对当前 shell
和子进程。需要长期一致的环境应写入项目包装脚本或 systemd
单元，并经过部署评审。不要全局修改
\passthrough{\lstinline!/etc/profile!} 来让所有用户继承未经验证的 CANN
和 Python 路径。

\section{14.
板端设备、服务、日志、测试与证据边界}\label{ux677fux7aefux8bbeux5907ux670dux52a1ux65e5ux5fd7ux6d4bux8bd5ux4e0eux8bc1ux636eux8fb9ux754c}

\subsection{系统状态报告}\label{ux7cfbux7edfux72b6ux6001ux62a5ux544a}

\begin{lstlisting}[language=bash]
uname -a
cat /etc/os-release
hostname
date --iso-8601=seconds
free -h
swapon --show
df -h
uptime
ps -p 1 -o pid,comm,args
npu-smi info
\end{lstlisting}

这份状态报告用于记录测试环境，不是性能测试。把它和后续
ATC、ACL、精度和性能命令放在同一报告目录中，便于追溯。

\subsection{摄像头和 V4L2}\label{ux6444ux50cfux5934ux548c-v4l2}

\begin{lstlisting}[language=bash]
ls -l /dev/video*
groups
v4l2-ctl --device=/dev/video0 --list-formats-ext
v4l2-ctl --device=/dev/video0 --all
\end{lstlisting}

\passthrough{\lstinline!/dev/video0!}
不一定是目标摄像头，设备编号可能随枚举顺序变化。先记录实际节点、格式、分辨率和帧率，再运行采集程序。\passthrough{\lstinline!--source 0!}
的含义由 V4L2 枚举顺序决定，不能仅凭一次运行推断为固定硬件。

无权限时检查组：

\begin{lstlisting}[language=bash]
groups
sudo usermod -aG video "$USER"
\end{lstlisting}

重新登录后再验证。不要用 root 启动整个图形程序或 Python
服务来绕开设备权限。

\subsection{USB 和音频设备}\label{usb-ux548cux97f3ux9891ux8bbeux5907}

\begin{lstlisting}[language=bash]
lsusb
lsusb -t
cat /proc/asound/cards
aplay -l
arecord -l
ls -l /dev/snd
\end{lstlisting}

音频设备编号和名称必须以当前板端输出为准。\passthrough{\lstinline!aplay -l!}
只列出播放设备，\passthrough{\lstinline!arecord -l!}
列出采集设备。播放测试会改变音量或产生声音，执行前确认扬声器、耳机和实验环境。

USB 采集设备通常需要独占。确认设备节点后查看占用进程：

\begin{lstlisting}[language=bash]
lsusb
usb_node="/dev/bus/usb/001/002"
test -e "$usb_node"
fuser -v "$usb_node"
\end{lstlisting}

\passthrough{\lstinline!fuser -v!}
只用于定位占用者，确认进程属于哪个程序后再处理。不要按模糊名称批量结束进程。

\subsection{服务和端口}\label{ux670dux52a1ux548cux7aefux53e3}

\begin{lstlisting}[language=bash]
systemctl status ssh --no-pager
systemctl --failed
systemctl list-units --type=service --state=running
ss -ltnp
ss -lunp
\end{lstlisting}

\passthrough{\lstinline!systemctl --failed!}
显示失败单元，\passthrough{\lstinline!ss -ltnp!} 显示 TCP
监听及进程，\passthrough{\lstinline!ss -lunp!} 显示 UDP 监听。服务绑定到
\passthrough{\lstinline!127.0.0.1!} 时只接受本机连接，绑定到
\passthrough{\lstinline!0.0.0.0!}
则对所有网卡开放。只有可信实验网络中的明确需求才使用
\passthrough{\lstinline!0.0.0.0!}。

\subsection{进程和日志}\label{ux8fdbux7a0bux548cux65e5ux5fd7}

\begin{lstlisting}[language=bash]
pgrep -af 'uvicorn|python|docker|npu'
ps -ef | grep -E 'python|uvicorn|docker' | grep -v grep
journalctl -b -p err --no-pager | tail -n 100
journalctl -u ascend310-demo.service -b --no-pager | tail -n 80
sudo dmesg -T | tail -n 200
\end{lstlisting}

\passthrough{\lstinline!journalctl -b!}
查看本次启动，\passthrough{\lstinline!-u!} 查看指定 systemd
单元，\passthrough{\lstinline!-p err!}
过滤错误级别。\passthrough{\lstinline!dmesg!} 查看内核消息，涉及
USB、摄像头、网络和驱动问题时尤其有用。公开报告只保留与故障相关的行，并删除用户名、IP、token、真实图像路径等隐私信息。

\subsection{安全停止服务}\label{ux5b89ux5168ux505cux6b62ux670dux52a1}

只停止已经确认的 PID：

\begin{lstlisting}[language=bash]
service_pid=12345
case "$service_pid" in
  ''|*[!0-9]*) printf 'invalid PID\n' >&2; exit 2 ;;
esac
test "$service_pid" -gt 1
ps -fp "$service_pid"
readlink -f "/proc/$service_pid/cwd"
kill -TERM "$service_pid"
for _ in $(seq 1 20); do
  kill -0 "$service_pid" 2>/dev/null || break
  sleep 1
done
\end{lstlisting}

先确认 PID 对应的命令、工作目录和用户，再发送
\passthrough{\lstinline!TERM!}。不要执行
\passthrough{\lstinline!pkill python!}、\passthrough{\lstinline!killall python!}
或按模糊名称批量结束进程，板端可能同时运行多个实验。

\subsection{本地测试}\label{ux672cux5730ux6d4bux8bd5}

开发机可以执行语法检查和纯 Python 测试：

\begin{lstlisting}[language=bash]
python -m py_compile app.py
python -m compileall -q package_or_directory
python -m pytest -q
python -m unittest discover -s tests -v
git diff --check
\end{lstlisting}

前端项目按锁文件执行：

\begin{lstlisting}[language=bash]
npm ci
npm test
npm run build
npm run test:e2e
\end{lstlisting}

真实板端硬件测试必须有显式开关，且只在板端执行。例如项目定义了硬件测试变量时：

\begin{lstlisting}[language=bash]
BOARD_HARDWARE_TESTS=1 python -m pytest -q tests/test_board_hardware.py
\end{lstlisting}

变量名应以项目实际脚本为准。普通单元测试不应隐式访问
NPU、摄像头、串口、音频设备或真实个人数据。

\subsection{证据门}\label{ux8bc1ux636eux95e8}

{ % do not increment counter
\begin{longtable}[]{@{}
  >{\raggedright\arraybackslash}p{(\linewidth - 4\tabcolsep) * \real{0.3333}}
  >{\raggedright\arraybackslash}p{(\linewidth - 4\tabcolsep) * \real{0.3333}}
  >{\raggedright\arraybackslash}p{(\linewidth - 4\tabcolsep) * \real{0.3333}}@{}}
\toprule\noalign{}
\begin{minipage}[b]{\linewidth}\raggedright
证据门
\end{minipage} & \begin{minipage}[b]{\linewidth}\raggedright
能证明
\end{minipage} & \begin{minipage}[b]{\linewidth}\raggedright
不能证明
\end{minipage} \\
\midrule\noalign{}
\endhead
\bottomrule\noalign{}
\endlastfoot
Python 语法检查 & 文件可解析 & 运行时依赖、板端设备、模型正确性 \\
HTTP 200 & 路由和服务进程可达 & NPU 推理成功、精度和性能 \\
\passthrough{\lstinline!npu-smi info!} & 设备可被管理工具读取 & 某个
OM、模型或应用程序已经验证 \\
\passthrough{\lstinline!import acl!} & PyACL 可在当前解释器导入 & ACL
context、具体模型和算子链通过 \\
ATC 退出码为 0 & 转换流程返回成功 & OM 数值和任务精度正确 \\
OM 摘要一致 & 文件字节未变化 & 模型可加载或输入输出合同正确 \\
ACL 数值烟测 & 指定输入的数值关系在容差内 & 数据集精度、延迟和吞吐 \\
任务精度测试 & 指定数据集和协议上的指标 & 其他数据集和硬件层级的结果 \\
性能测试 & 指定参数下的延迟或吞吐 & 精度和端到端业务正确性 \\
UI 烟测 & 页面或控件可以交互 & 后端真实硬件推理已经完成 \\
\end{longtable}
}

每份报告至少记录硬件型号和算力层级、CANN 版本、模型 ID
和精度、输入合同、预热次数、循环次数、重复次数、百分位方法、输入来源、输出路径和原始日志。

\section{15.
故障排查与禁止命令模式}\label{ux6545ux969cux6392ux67e5ux4e0eux7981ux6b62ux547dux4ee4ux6a21ux5f0f}

\subsection{常见故障}\label{ux5e38ux89c1ux6545ux969c}

{ % do not increment counter
\begin{longtable}[]{@{}
  >{\raggedright\arraybackslash}p{(\linewidth - 4\tabcolsep) * \real{0.3333}}
  >{\raggedright\arraybackslash}p{(\linewidth - 4\tabcolsep) * \real{0.3333}}
  >{\raggedright\arraybackslash}p{(\linewidth - 4\tabcolsep) * \real{0.3333}}@{}}
\toprule\noalign{}
\begin{minipage}[b]{\linewidth}\raggedright
现象
\end{minipage} & \begin{minipage}[b]{\linewidth}\raggedright
先执行
\end{minipage} & \begin{minipage}[b]{\linewidth}\raggedright
判断重点
\end{minipage} \\
\midrule\noalign{}
\endhead
\bottomrule\noalign{}
\endlastfoot
SSH 无法连接 &
\passthrough{\lstinline!ping!}、\passthrough{\lstinline!ip addr!}、\passthrough{\lstinline!ip route!}、\passthrough{\lstinline!systemctl status ssh!}
& 是否同网段、服务是否启动、22 端口是否监听 \\
known\_hosts 冲突 & \passthrough{\lstinline!ssh-keygen -R <host>!} &
只删除目标主机记录，并核对新指纹 \\
VNC 无法连接 & 检查共享开关、板端 IP、VNC 服务和显示输出 & VNC 与 SSH
的端口和权限边界不同 \\
VNC 分辨率低 &
\passthrough{\lstinline!xrandr!}、\passthrough{\lstinline!xrandr --fb 1024x768!}
& 是否超过显示能力，是否影响画面质量 \\
WiFi 切换后失联 & 检查热点、网卡、连接优先级和新 IP &
一个网卡通常不能同时做热点和连接上游 WiFi \\
无法 ping 通板端 &
\passthrough{\lstinline!ip addr!}、\passthrough{\lstinline!ip route!}、\passthrough{\lstinline!ping!}
& 同网段、路由、热点和地址冲突 \\
NFS 挂载失败 &
\passthrough{\lstinline!ping!}、\passthrough{\lstinline!showmount!}、\passthrough{\lstinline!exportfs -v!}、检查路径
& 导出范围、目录路径、客户端权限和失效挂载 \\
NFS 服务端关机 & 先 \passthrough{\lstinline!umount!} 再重新
\passthrough{\lstinline!mount!} & 失效挂载不能直接覆盖 \\
Docker 容器不在运行 &
\passthrough{\lstinline!docker ps -a!}、\passthrough{\lstinline!docker logs!}
& 容器状态、退出码和重启策略 \\
容器中无图形窗口 & 检查 \passthrough{\lstinline!DISPLAY!}、X11
认证、容器环境变量 & 图形显示与业务推理是两个问题 \\
\passthrough{\lstinline!import acl!} 失败 &
\passthrough{\lstinline!command -v python!}、\passthrough{\lstinline!source set\_env.sh!}、重新导入
& 解释器和 CANN 是否来自同一环境 \\
ATC 找不到算子 & \passthrough{\lstinline!atc --version!}、检查
framework、soc\_version、输入形状和日志 &
保存失败命令和模型合同，不生成伪 OM \\
\passthrough{\lstinline!npu-smi!} 显示 Alarm &
保存状态报告并检查具体错误 & 告警是诊断背景，不能单独作为测试结论 \\
摄像头无图像 &
\passthrough{\lstinline!/dev/video*!}、\passthrough{\lstinline!v4l2-ctl!}、\passthrough{\lstinline!groups!}
& 设备节点、格式、权限和实际尺寸 \\
USB 仪器被占用 &
\passthrough{\lstinline!lsusb!}、\passthrough{\lstinline!fuser -v <device>!}
& 确认占用者，不用宽泛 kill \\
音频无声音 &
\passthrough{\lstinline!aplay -l!}、\passthrough{\lstinline!pactl list short sinks!}
& 设备、默认 sink 和配置 profile \\
端口被占用 &
\passthrough{\lstinline!ss -ltnp!}、\passthrough{\lstinline!ps -fp <pid>!}
& 只停止确认属于目标服务的 PID \\
磁盘不足 &
\passthrough{\lstinline!df -h!}、\passthrough{\lstinline!du -sh <dir>!}
& 先归档报告，再清理隔离临时目录 \\
\end{longtable}
}

\subsection{禁止或必须审查的命令模式}\label{ux7981ux6b62ux6216ux5fc5ux987bux5ba1ux67e5ux7684ux547dux4ee4ux6a21ux5f0f}

以下命令模式不应未经确认直接执行：

\begin{itemize}
\tightlist
\item
  \passthrough{\lstinline!rm -rf!}
  指向家目录、仓库根目录、变量未验证的路径或正在运行的部署目录。
\item
  \passthrough{\lstinline!rsync --delete!}
  同步到未确认的远端目录，它可能删除模型、数据和报告。
\item
  \passthrough{\lstinline!pkill python!}、\passthrough{\lstinline!killall!}
  或模糊
  \passthrough{\lstinline!kill!}，它们可能终止其他案例和系统服务。
\item
  \passthrough{\lstinline!sudo pip install!}、\passthrough{\lstinline!sudo python!}，它们会绕过
  Conda 并改变系统包所有权。
\item
  \passthrough{\lstinline!chmod -R 777!} 和未核实所有者的递归
  \passthrough{\lstinline!chown!}，它们会放宽整个目录树的权限。
\item
  在开发机运行 CANN、ATC、PyACL、OM、摄像头或 NPU
  命令，并把失败归因于代码。
\item
  在一个终端加载 CANN 或 Conda，却在另一个未加载环境的终端启动服务。
\item
  启动 \passthrough{\lstinline!0.0.0.0!}、公网隧道、未鉴权网关或宽泛 X11
  授权后再处理个人数据。
\item
  使用 \passthrough{\lstinline!xhost +!}
  允许所有客户端访问显示服务器，而不是限制到当前用户。
\item
  直接使用 \passthrough{\lstinline!--privileged!}
  给容器开放全部宿主能力，而不是逐项授权设备。
\item
  把 \passthrough{\lstinline!npu-smi!} 摘要、HTTP 200、软件 tone、CPU
  回退或界面可达写成完整硬件验收。
\item
  未确认路径就复制真实照片、生物特征模板、数据库、模型、密钥或 token 到
  Git、截图和公开报告。
\end{itemize}

安全的操作顺序是：确认机器和环境，解析绝对路径，预览输入与目标，执行最小范围命令，保存原始输出，再把经过审查的结论写入报告。这样才能让
Ubuntu 命令成为可复现的操作步骤，而不是无法追溯的部署记录。
